\chapter*{Abstract}
\thispagestyle{empty}
RNA molecules play essential biological roles whose functions are critically determined by their three-dimensional (3D) structures. However, experimental determination of RNA 3D structures remains challenging, and existing computational approaches still struggle to achieve accurate, generalizable predictions across diverse RNA families and lengths. In particular, deterministic deep learning models often suffer from limited robustness to conformational heterogeneity, length scaling, and RNA-specific geometric constraints.

In this thesis, we propose Diffold, a diffusion-based generative framework for RNA 3D structure prediction. Diffold integrates (i) the evolutionary and geometric representations produced by the RhoFold+ backbone, (ii) structural priors derived from the large-scale RNA language model RNA-FM, and (iii) an Elucidated Diffusion Model (EDM) formulation to generate all-atom RNA coordinates through iterative denoising. A hierarchical token–atom denoising architecture is employed to simultaneously capture long-range residue-level dependencies and local atomic geometry, enabling physically plausible structure generation in Cartesian coordinate space.

Diffold is evaluated on the combined CASP16 RNA targets and the RNA-Benchmark dataset. Across all major metrics, Diffold consistently outperforms the strong baseline RhoFold+. Specifically, Diffold reduces RMSD from 3.85 $\text{\AA}$ to 2.29 $\text{\AA}$ and improves TM-score from 0.309 to 0.585, indicating substantially improved global fold recovery. Additional gains are observed in GDT-TS and lDDT, together with a marked reduction in steric clashes, reflecting enhanced local geometric consistency and physical plausibility. Further analyses demonstrate that Diffold exhibits weaker dependence on training-set structural similarity and improved robustness to increasing sequence length, suggesting stronger generalization than deterministic backbone-based methods.

Finally, we show that a lightweight AMBER-based relaxation step effectively corrects local covalent-geometry artifacts, such as elongated O3′–P bonds, without materially affecting global accuracy. Overall, these results indicate that diffusion-based generative modeling, when conditioned on rich RNA-specific representations, provides a promising and extensible strategy for accurate and generalizable RNA 3D structure prediction.